\documentclass[11pt,a4paper]{article}
\usepackage[utf8]{inputenc}\usepackage[T1]{fontenc}
\usepackage{amsmath,amssymb,amsthm,microtype,xcolor}
\usepackage[margin=2.5cm]{geometry}\linespread{1.06}\setlength{\emergencystretch}{3em}\hfuzz=2.5pt
\usepackage[colorlinks=true,linkcolor=blue!60!black,citecolor=blue!60!black]{hyperref}
\newtheorem{theorem}{Theorem}\newtheorem{proposition}[theorem]{Proposition}\newtheorem{lemma}[theorem]{Lemma}
\newtheorem{corollary}[theorem]{Corollary}\theoremstyle{remark}\newtheorem{remark}[theorem]{Remark}
\newcommand{\Z}{\mathbb{Z}}\newcommand{\R}{\mathbb{R}}\newcommand{\Off}{\mathrm{Off}}\newcommand{\e}{\mathrm{e}}
\newcommand{\todo}[1]{\textcolor{red!70!black}{\small\sffamily[#1]}}
\title{The pair singular series of a polynomial V:\\ the dilated pieces below $H^{1/2}$}
\author{Jasper Reichardt\thanks{Independent researcher. Correspondence: jasperreichardt@gmail.com.}}
\date{16 September 2026 --- working draft, not for circulation}
\begin{document}\maketitle
\begin{abstract}
\todo{Draft. The proofs of Theorems~\ref{thm:M} and~\ref{thm:S} are in research/explore/PAPER-V-thmM.md and PAPER-V-smallend-draft.md and are under adversarial reading;
nothing here is to be circulated before those readings are recorded.}
For a monic irreducible quadratic $f$ without fixed prime divisor, part~III decomposed the Ces\`aro off-diagonal $\Off^\ast_f(H)$ into pieces indexed by the dilation $u$ and
proved that the pieces with $u>H^{1/2+\varepsilon}$ contribute a main term $c_{\mathrm{off}}(f)H$. This paper treats dilations below that threshold. We prove that the pieces with
$H^{0.47}<u\le H^{0.53}$ contribute $O(H^{1-\delta})$, so that $\sum_{u>H^{0.47}}w(u)\mathcal P_u(H/u)=c_{\mathrm{off}}(f)H+O(H^{1-\delta})$; the input is a bilinear estimate for
Kloosterman fractions of Bourgain and Garaev together with a bound for the additive energy of reciprocals, valid for every modulus, which we deduce from an argument of
Cilleruelo and Garaev. For negative discriminants we prove a second result at the other end: the pieces with $u\le H^{1/3-c}$ contribute $O(H^{1-\delta(c)})$, by running the
automorphic kernel method of Grimmelt and Merikoski with the leading coefficient $u^2$ tracked. What remains is a neighbourhood of $u=H^{1/3}$.
\end{abstract}
\section{Introduction}
\todo{Working draft: the proofs are in research/explore/PAPER-V-thmM.md and PAPER-V-smallend-draft.md and are under adversarial reading.}

\paragraph{The problem.} Let $f=t^2+bt+c\in\Z[t]$ be monic, irreducible, without fixed prime divisor, of discriminant $D$. Part~III~\cite{ReichardtIII} wrote the Ces\`aro
off-diagonal of the pair singular series as a sum of \emph{pieces},
\[
\Off^\ast_f(H)=P(1)\sum_{u}\lambda(u)\omega(u)\,\mathcal P_u(H/u),
\]
one for each squarefree \emph{dilation} $u$, the piece $\mathcal P_u(Y)$ being the same problem at length $Y=H/u$ with the roots of $x^2\equiv D$ replaced by those of
$u^2x^2\equiv D$. Each piece is $O(Y)$, and the weights satisfy $\sum_{u\le H}|w(u)|/u\asymp\log H$: the trivial bound for the whole sum is $H\log H$, which is exactly the size
of the statement to be proved. Every piece must therefore be small, and the difficulty is uniformity in $u$.

\paragraph{Exponent coordinates.} Write $u=H^{\alpha}$. Each method available for the pieces gives a saving $H^{-s(\alpha)}$ over the trivial bound on an interval of $\alpha$,
and degenerates at the ends of that interval. In these coordinates part~III proves $s>0$ for $\alpha>2/3$ (averaging the dilations along $u$), and the present author's
extension of that theorem, recorded in the revised part~III, gives $s>0$ for $\alpha>1/2$ by a dispersion argument with Weil's bound. Below $1/2$ both mechanisms fail: the
incomplete Kloosterman sums that dispersion produces have modulus $[d_1,d_2]\asymp Y^2$, and Weil's bound is trivial once the range of $u$ is shorter than the square root of
that modulus. This paper supplies two further mechanisms, one on each side of $\alpha=1/3$.

\paragraph{Results.} Theorem~\ref{thm:M} treats $\alpha$ near $1/2$: the weight $w$ is decomposed as $1*\chi_D*\gamma$, and the resulting bilinear forms in $u$ are estimated by
Weil's bound when one factor is long (Type~I) or short (Type~II), and, in the remaining balanced case, by the bilinear bound of Bourgain and Garaev~\cite{BourgainGaraev2014}
for Kloosterman fractions. That bound needs an estimate for the additive energy of reciprocals; the one in~\cite{BourgainGaraev2014} is too weak at the sizes that occur here,
and Lemma~\ref{lem:E} supplies a stronger one, valid for every modulus, by transferring an argument of Cilleruelo and Garaev~\cite{CilleruelooGaraev2011} for points on modular
hyperbolas. Theorem~\ref{thm:S} treats small $\alpha$ for negative discriminants: piece $u$ is, after removing the content, a Type~I sum for the polynomial $u^2\ell^2+|D|$ in
the sense of Grimmelt and Merikoski~\cite{GrimmeltMerikoski2025quadratic}, and their automorphic kernel theorem~\cite{GrimmeltMerikoski2025kernel} applies at level $u^2e$ once
the two kernel quantities are estimated with the dilation tracked. The diagonal is the number of orbits of the family, $\asymp u$; the off-diagonal, which their general bound
would give as $\asymp u^2$, is in fact bounded independently of $u$ (Lemma~\ref{lem:Aoff}), because the distances inside the family are quantised in steps of $1/(4h)$ whatever
the dilation.

\paragraph{What remains.} Both mechanisms degenerate at $\alpha=1/3$, from opposite sides and for unrelated reasons: there the balanced Type~II atoms sit at $d^{1/4}$, where the
energy input is not strong enough, and the kernel method's diagonal $\asymp u$ meets the trivial bound. The final section explains what the residual band costs. Two points are
worth recording. A band of width $2\delta$ in $\alpha$ costs $\delta\cdot H(\log H)^{1-c}\log\log H$, not $\delta H\log H$, because part~III's Theorem~E applies piece by piece
inside it; so a failing method never makes the total worse. And if the savings on both sides vanish only \emph{at} $\alpha=1/3$, with implied constants that are powers of
$\log H$, then the blocks near the point contribute $O(H)$ in total, which would give $\Off^\ast_f(H)=O(H)$ --- Hypothesis~(E) of part~I~\cite{ReichardtI}, a bounded
second-order term in the Ces\`aro asymptotic. The constant $A_f$ needs strictly more: a saving at the point itself.

\section{Notation}
We follow part~III~\cite{ReichardtIII}. For a prime $p$ let $\omega(p)$ be the number of roots of $f$ modulo $p$, $\lambda(p)=p/(p-2\omega(p))$, $P_p=1-\omega(p)^2/(p-\omega(p))^2$,
$P(1)=\prod_pP_p$ and $w(u)=P(1)\lambda(u)\omega(u)$; $w$ is supported on squarefree $u$ and $\sum_{u\le t}|w(u)|\ll t$. A modulus $d'$ is \emph{admissible} if it is squarefree,
all its prime factors are split, and $(d',2Du)=1$; $R^{(u)}_{d'}=\{x\bmod d':u^2x^2\equiv D\}$, $\rho_k(d';a)=\sum_{r\in R_{d'}}\e(kar/d')$, and $\langle x\rangle_{d'}\in[1,d']$ is
the representative of $x$. With $B^{(Y)}_{d'}(m)=\sum_{h\le Y,\,h\equiv m\,(d')}(Y-h)-Y^2/(2d')$, the piece is
\[
\mathcal P_u(Y)=\sum_{d'>1\ \mathrm{admissible}}\frac{\lambda(d')}{d'}\sum_{x\in R^{(u)}_{d'}}B^{(Y)}_{d'}(\langle x\rangle_{d'}),\qquad Y=H/u .
\]
We write $G_{d'}(a,Y)=\sum_{r\in R_{d'}}B^{(Y)}_{d'}(\langle ar\rangle_{d'})$, $\bar G_{d'}(Y)$ for its mean over the units $a$, $G^\circ=G-\bar G$, and $\xi_k(d',N)=(1-\cos(2\pi
kN/d'))/\sin^2(\pi k/d')$ for the coefficients of the finite Fourier expansion of $B^{(N)}$ (part~III, Lemma~finfourier). Throughout, $x^\ast$ denotes the inverse of $x$ modulo the
modulus at hand, $e_m(z)=\e(z/m)$, and $\varepsilon>0$ is arbitrarily small. Parameters: $\delta=10^{-3}$, $\eta=\delta/10$, $\delta_3=\delta$, $\varepsilon'=\eta/7$,
$\kappa=\delta/2$, $K_1=H^{2\eta}$.
\section{The energy of reciprocals and bilinear Kloosterman fractions}
\begin{lemma}[energy of reciprocals, every squarefree modulus]\label{lem:E}
Let $m\ge2$ be squarefree, $1\le N\le m$, and $J_4(N)=\#\{x\in[1,N]^4:(x_i,m)=1,\ x_1^\ast+x_2^\ast\equiv x_3^\ast+x_4^\ast\ (\mathrm{mod}\ m)\}$. Then for every $\varepsilon>0$,
\[
J_4(N)\ll_\varepsilon m^{\varepsilon}\bigl(N^2+N^{7/2}m^{-1/2}\bigr).
\]
\end{lemma}
\todo{Proof written in PAPER-V-thmM.md (Cilleruelo--Garaev's argument, transferred to composite moduli); two readings recorded in PROOFS-uniform.md §28, §30.}

\begin{corollary}\label{cor:B}
For squarefree $m$, $(a,m)=1$, $m^{1/3}\le N_1,N_2\le m$ and coefficients $|\alpha_i|\le1$ supported on units,
$\sum_{x_1\le N_1}\sum_{x_2\le N_2}\alpha_1(x_1)\alpha_2(x_2)e_m(ax_1^\ast x_2^\ast)\ll_\varepsilon m^{\varepsilon}(N_1N_2)^{15/16}$.
\end{corollary}
\todo{From \cite{BourgainGaraev2014}, proof of Theorem~3, together with Lemma~\ref{lem:E}.}
\section{The pieces around $u=H^{1/2}$}
\begin{theorem}\label{thm:M}
Let $f$ be as above. There is $\delta_M>0$ (any $\delta_M<4\eta/7$ with the parameters above, so $\delta_M=5\cdot10^{-5}$ is admissible) such that
\[
\sum_{H^{0.47}<u\le H^{0.53}}w(u)\,\mathcal P_u(H/u)\ \ll_f\ H^{1-\delta_M}.
\]
Consequently, with part~III's Theorem on the range $u>H^{1/2+\varepsilon}$ applied at $\varepsilon=0.03$,
$\sum_{u>H^{0.47}}w(u)\mathcal P_u(H/u)=c_{\mathrm{off}}(f)H+O(H^{1-\delta_M})$.
\end{theorem}
\todo{Full proof in research/explore/PAPER-V-thmM.md (with REVISION 1 and REVISION 2 merged); three independent readings, recorded in PROOFS-uniform.md §28, §30, §35, the last
confirming the merged text. What follows is the skeleton to be expanded into the paper text.}

\paragraph{Structure of the proof.} Fix a smooth dyadic partition of the range of $u$, with ramps of relative width $H^{-\kappa}$ at the two ends; the transition zones are
$\ll H^{1-\kappa}(\log H)^C$ by the trivial bound $\mathcal P_u(Y)\ll Y$. For one block $U\in[H^{0.47}/4,2H^{0.53}]$, $Y=H/U$, $a=\log U/\log H\in[0.46,0.54]$:
\begin{enumerate}
\item[R1.] \emph{Reduction.} The dilation average contributes $\ll U(\log H)^C$; the moduli $d'>YH^{\eta}$ are part~III's Step~3, $\ll H^{1-\eta+3\varepsilon'}$; the moduli
$d'\le Y^{1-\delta}$ are $\ll UY^{1-\delta}$ by Shiu's theorem; on the remaining range the finite Fourier expansion applies, the frequencies $k>K_1$ cost $H^{1-\eta}$, and a
Mellin transform separates the coefficients $\xi_k(d',H/u)$ from $u$, at a cost of $\ll Y(\log H)^3$ after summing over $k$ and $d'$.
\item[R5.] \emph{Decomposition.} $w=c_f(1*\chi_D*\gamma)$ with $\sum_n|\gamma(n)|n^{-1/2-\varepsilon}<\infty$; the factor $n_3>H^{\delta_3}$ is removed trivially; the atoms
$n_1$ (weight $1$) and $n_2$ (weight $\chi_D$) are split into smooth dyadic boxes. Three cases exhaust the boxes.
\item[R2.] \emph{Type I:} an atom of length $N\ge Y^{1/2}H^{3\delta}$ gives, by Lemma~K applied for each modulus and cofactor, a saving $H^{-1.79\delta}$.
\item[R3.] \emph{Type II:} an atom of length $U_2\in[H^{3\delta},UY^{-1/2}H^{-3\delta}]$ gives, by Cauchy--Schwarz over $(d',u_1)$ and Weil's bound for the resulting incomplete
Kloosterman sums, a saving $H^{-1.45\delta}$.
\item[R4.] \emph{The residual:} both atoms of size between $UY^{-1/2}H^{-3\delta}$ and $Y^{1/2}H^{3\delta}$. Then both exceed $m^{1/3}$ for $m=d'/(k,d')$, and
Corollary~\ref{cor:B} gives a saving $H^{-0.028}$.
\end{enumerate}
Summing the three savings against the $O(\log H)$ blocks, the binding constraint is part~III's Step~3, which yields $\delta_M<4\eta/7$.
\section{The small dilations}
\begin{lemma}[the off-diagonal at one level does not grow with the dilation]\label{lem:Aoff}
Let $h\ge1$, $u$ squarefree with $(u,2h)=1$, $a=u^2$, and let $F_u$ be the set of positive definite forms $g=\bigl(\begin{smallmatrix}A&B\\B&C\end{smallmatrix}\bigr)$ with
$AC-B^2=ah$, $a\mid B$, $a\mid C$. For $Z\ge1$,
\[
\sum_{g_2\in\Gamma_0(a)\backslash F_u}\ \sum_{\substack{g_1\in F_u,\ g_1\ne g_2\\ u(g_1,g_2)\le Z}}\bigl(1+u(g_1,g_2)\bigr)^{-1/2}\ \ll_\varepsilon\ (hZu)^{\varepsilon}h^{O(1)}\bigl(1+Z^{1/2}\bigr),
\]
uniformly in $u$, where $u(\cdot,\cdot)$ is the point-pair invariant.
\end{lemma}
\todo{Proof in PAPER-V-smallend-draft.md (undilation to $\Gamma_0^0(u)$, lift to level one, a local count of common orthogonal lines, and a sum over the quantised distances);
under reading.}

\begin{theorem}\label{thm:S}
Let $f$ be as above with $D<0$. For every $c>0$ there is $\delta_S(c)>0$ with
$\sum_{u\le H^{1/3-c}}w(u)\mathcal P_u(H/u)\ll_f H^{1-\delta_S(c)}$.
\end{theorem}
\todo{Proof outline S1--S8 in PAPER-V-smallend-draft.md, resting on \cite{GrimmeltMerikoski2025kernel} Theorem~8.1 and \cite{GrimmeltMerikoski2025quadratic} Lemma~3.1, with
Lemma~\ref{lem:Aoff} and the diagonal count; $\delta_S(c)\asymp c$ once the exponent in the kernel theorem's $\delta^{-O(1)}$ is named. Under reading.}
\section{What remains: the band at $u=H^{1/3}$}
Theorems~\ref{thm:M} and~\ref{thm:S} leave a neighbourhood of $\alpha=1/3$ uncovered, and the two mechanisms fail there for unrelated reasons. On the upper side the balanced
Type~II atoms sit at $d^{1/4}$, where the available energy input for reciprocals is exactly at its threshold: what would be needed is a bound for the six-fold energy
$J_6(N)$ below $N^4$ at $N\asymp m^{1/4}$, which is not known even for prime moduli. On the lower side the kernel method's diagonal, the number of orbits of the family, is
$\asymp u$, and at $u\asymp Y^{1/2}$ it meets the trivial bound.

Three remarks fix what such a band costs, and what closing it would buy.

\begin{remark}[a band is cheap, and a failing method is free]
Inside any range of $\alpha$ one still has part~III's Theorem~E piece by piece, $\mathcal P_u(H/u)\ll(H/u)(\log H)^{-c}\log\log H$ uniformly in $u\le H^{1-\varepsilon}$. Hence a band
$[H^{\alpha_0-\delta},H^{\alpha_0+\delta}]$ on which nothing better is available costs $\ll_f\delta\cdot H(\log H)^{1-c}\log\log H$, not $\delta H\log H$: a mechanism that
degenerates can never make the total worse than what is already proved, and any band whose width tends to $0$ improves Theorem~E by that factor.
\end{remark}

\begin{remark}[a point gap costs $O(H)$]
Suppose the savings on the two sides vanish only at $\alpha=1/3$, say like $H^{-c|\alpha-1/3|}$, with implied constants that are powers of $\log H$. Summing over dyadic blocks,
$\sum_kH\,\mathrm{e}^{-c|k\log2-\frac13\log H|}=H\bigl(1+2/(\mathrm{e}^{c\log2}-1)\bigr)=O(H/c)$: the blocks near the point contribute $O(H)$ in total. Since the main term
$c_{\mathrm{off}}(f)H$ comes from $u>H^{1/2+\varepsilon}$ and the remaining blocks are error terms, this would give $\Off^\ast_f(H)=O(H)$, which is Hypothesis~(E) of
part~I~\cite{ReichardtI} and yields $\sum_{h\le H}(1-h/H)(S_f(h)-C(f)^2)=-\tfrac12C(f)\log H+O_f(1)$: a bounded second-order term.
\end{remark}

\begin{remark}[what $A_f$ needs]
The constant of Conjecture~1 of part~I requires $\Off^\ast_f(H)=c_{\mathrm{off}}(f)H+o(H)$, hence a saving factor $\gg\log H$ uniformly on $|\alpha-1/3|\ll\log\log H/\log H$; a
fixed positive power saving at $\alpha=1/3$ itself suffices. A linearly vanishing saving, however small the rate, gives $O(H)$ but never $o(H)$.
\end{remark}

\begin{remark}[what is actually binding]
An audit of every implied constant in the chain (research/explore/PROOFS-uniform.md §36) shows that no divisor-size loss constrains the argument: the divisor factors all occur
under sums and average to powers of $\log H$, the kernel theorem's $\delta^{-O(1)}$ is $\delta^{-20}$, its $(AD)^{o(1)}$ is $(\log AD)^{O(1)}$, and its $q^{o(1)}$ is $1$ in the
present application. What binds is the constant in Bourgain and Garaev's H\"older step, which grows as $k_1\asymp 2/(9|\alpha-1/3|)\to\infty$: the band on which the estimate is
effective has width $\asymp(\log\log H/\log H)^{1/3}$. With the in-band fallback on part~III's Theorem~E this gives, once the map is written out,
$\Off^\ast_f(H)\ll H(\log H)^{0.63+o(1)}$ for $D<0$, while everything outside the band sums to $O(H)$. The step from there to $\Off^\ast_f(H)=O(H)$ therefore needs one thing only:
an estimate that survives at the exceptional sizes $N\asymp m^{1/4}$ --- an energy bound $J_6$ below $N^4$ for composite moduli, or a different tool there. A level-averaged
spectral bound for the twist $\mathrm{e}(cr\bar d/x)$ with $x\asymp Y^{1/2}$ would serve the same purpose; neither is in the literature.
\end{remark}

\bibliographystyle{plain}\bibliography{refs}
\end{document}
